%\usepackage{lmodern}
%\usepackage[T1]{fontenc}
%\usepackage{concrete}
%\usepackage{eulervm}
%\usefonttheme[onlymath]{serif}
%\usefonttheme{professionalfonts}
%\usecolortheme[named=black]{structure}
%\definecolor{red}{rgb}{0,0,0}
%\definecolor{blue}{rgb}{0,0,0}


\documentclass[notes=show,compress,beamer]{beamer}
%%%%%%%%%%%%%%%%%%%%%%%%%%%%%%%%%%%%%%%%%%%%%%%%%%%%%%%%%%%%%%%%%%%%%%%%%%%%%%%%%%%%%%%%%%%%%%%%%%%%%%%%%%%%%%%%%%%%%%%%%%%%%%%%%%%%%%%%%%%%%%%%%%%%%%%%%%%%%%%%%%%%%%%%%%%%%%%%%%%%%%%%%%%%%%%%%%%%%%%%%%%%%%%%%%%%%%%%%%%%%%%%%%%%%%%%%%%%%%%%%%%%%%%%%%%%
\usepackage{amsmath}
\usepackage{amsfonts}
\usepackage{mathpazo}
\usepackage{hyperref}
\usepackage{multimedia}
\usepackage{harvard}
\usepackage{color}
\usepackage{bibentry}

\setcounter{MaxMatrixCols}{10}
%TCIDATA{OutputFilter=LATEX.DLL}
%TCIDATA{Version=5.50.0.2953}
%TCIDATA{<META NAME="SaveForMode" CONTENT="1">}
%TCIDATA{BibliographyScheme=BibTeX}
%TCIDATA{Created=Tuesday, December 04, 2007 18:43:18}
%TCIDATA{LastRevised=Monday, May 03, 2010 13:01:25}
%TCIDATA{<META NAME="GraphicsSave" CONTENT="32">}
%TCIDATA{<META NAME="DocumentShell" CONTENT="Other Documents\SW\Slides - Beamer">}
%TCIDATA{Language=American English}
%TCIDATA{CSTFile=beamer.cst}

\newenvironment{stepenumerate}{\begin{enumerate}[<+->]}{\end{enumerate}}
\newenvironment{stepitemize}{\begin{itemize}[<+->]}{\end{itemize} }
\newenvironment{stepenumeratewithalert}{\begin{enumerate}[<+-| alert@+>]}{\end{enumerate}}
\newenvironment{stepitemizewithalert}{\begin{itemize}[<+-| alert@+>]}{\end{itemize} }
\useoutertheme{shadow}
\usetheme{boxes}
\addheadbox{structure}{ }
\addfootbox{structure}{ }
\addfootbox{structure}{ }
\addfootbox{structure}{  \hspace{1in} \color{gray} \scriptsize \insertframenumber /\inserttotalframenumber}
\useoutertheme[footline=empty,subsection=false]{miniframes}
\addtobeamertemplate{headline}
{}
{
\vskip1pt
}
\setbeamertemplate{navigation symbols}{}
\input{tcilatex}
\let\cite=\citeasnoun
\hypersetup{pdfpagemode=UseNone}

\begin{document}


%TCIMACRO{%
%\TeXButton{TITLE}{\title{
%Economics 326 \\
%Methods of Empirical Research in Economics\\
%\bigskip
%Lecture 14: Hypothesis testing in the multiple regression model, Part 2
%}}}%
%BeginExpansion
\title{
Economics 326 \\
Methods of Empirical Research in Economics\\
\bigskip
Lecture 14: Hypothesis testing in the multiple regression model, Part 2
}%
%EndExpansion

%TCIMACRO{%
%\TeXButton{AUTHORS AND AFFILIATION}{\author{Vadim Marmer\\
%University of British Columbia
%}}}%
%BeginExpansion
\author{Vadim Marmer\\
University of British Columbia
}%
%EndExpansion

%TCIMACRO{%
%\TeXButton{CREATES THE TITLE PAGE}{\setbeamertemplate{blocks}[rounded][shadow=true]
%\setbeamercolor{titlepg}{fg=white,bg=blue!50!gray}
%\setbeamercolor{title}{fg=white,bg=blue!50!gray}
%\setbeamerfont{title}{size=\Large}
%\setbeamerfont{inst}{size=\small}
%
%\frame[plain]{
%
%\begin{beamercolorbox}[center,shadow=true,rounded=true,sep=.05cm]{titlepg}
%\usebeamerfont{title}
%\center
%\inserttitle
%\end{beamercolorbox}
%
%\bigskip
%\bigskip
%\center{\insertauthor}
%\center{\insertdate}
%
%}}}%
%BeginExpansion
\setbeamertemplate{blocks}[rounded][shadow=true]
\setbeamercolor{titlepg}{fg=white,bg=blue!50!gray}
\setbeamercolor{title}{fg=white,bg=blue!50!gray}
\setbeamerfont{title}{size=\Large}
\setbeamerfont{inst}{size=\small}

\frame[plain]{

\begin{beamercolorbox}[center,shadow=true,rounded=true,sep=.05cm]{titlepg}
\usebeamerfont{title}
\center
\inserttitle
\end{beamercolorbox}

\bigskip
\bigskip
\center{\insertauthor}
\center{\insertdate}

}%
%EndExpansion

%TCIMACRO{%
%\TeXButton{\addtocounter{framenumber}{-1}}{\addtocounter{framenumber}{-1}}}%
%BeginExpansion
\addtocounter{framenumber}{-1}%
%EndExpansion

%TCIMACRO{\TeXButton{BeginFrame}{\begin{frame}[shrink]}}%
%BeginExpansion
\begin{frame}[shrink]%
%EndExpansion

\QTR{frametitle}{Multiple restrictions}

\begin{itemize}
\item Consider the model:%
\begin{multline*}
\ln \left( Wage_{i}\right) =\beta _{0}+\beta _{1}Experience_{i}+\beta
_{2}Experience_{i}^{2}+ \\
+\beta _{3}PrevExperience_{i}+\beta _{4}PrevExperience_{i}^{2}+\beta
_{5}Education_{i}+U_{i},
\end{multline*}%
where $Experience$ is the experience at current job, and $PrevExperience$ is
the previous experience.%
%TCIMACRO{\TeXButton{pause}{\pause}}%
%BeginExpansion
\pause%
%EndExpansion

\item Suppose that we want to test the null hypothesis that, after
controlling for the experience at current job and education, the previous
experience has no effect on wage:%
%TCIMACRO{\TeXButton{pause}{\pause}}%
%BeginExpansion
\pause%
%EndExpansion
\begin{equation*}
H_{0}:\beta _{3}=0,\beta _{4}=0.
\end{equation*}%
%TCIMACRO{\TeXButton{pause}{\pause}}%
%BeginExpansion
\pause%
%EndExpansion

\item We have 
%TCIMACRO{\TeXButton{\color{blue} CTRL+cbl}{\color{blue}}}%
%BeginExpansion
\color{blue}%
%EndExpansion
two 
%TCIMACRO{\TeXButton{\color{black} CTRL+cb}{\color{black}} }%
%BeginExpansion
\color{black}
%EndExpansion
restrictions on the model parameters.%
%TCIMACRO{\TeXButton{pause}{\pause}}%
%BeginExpansion
\pause%
%EndExpansion

\item The alternative hypothesis is that at least one of the coefficients, $%
\beta _{3}$ or $\beta _{4},$ is different from zero%
%TCIMACRO{\TeXButton{pause}{\pause}}%
%BeginExpansion
\pause%
%EndExpansion
:%
\begin{equation*}
H_{1}:\beta _{3}\neq 0\text{ or }\beta _{4}\neq 0.
\end{equation*}
\end{itemize}

%TCIMACRO{\TeXButton{EndFrame}{\end{frame}}}%
%BeginExpansion
\end{frame}%
%EndExpansion

%TCIMACRO{\TeXButton{BeginFrame}{\begin{frame}}}%
%BeginExpansion
\begin{frame}%
%EndExpansion

\QTR{frametitle}{%
%TCIMACRO{\TeXButton{$t$}{$t$}}%
%BeginExpansion
$t$%
%EndExpansion
-statistics and multiple restrictions}

\begin{itemize}
\item Let $T_{3}$ and $T_{4}$ be the $t$-statistics associated with the
coefficients of $PrevExperience$ and $PrevExperience^{2}:$%
\begin{equation*}
T_{3}=\frac{\hat{\beta}_{3}}{SE\left( \hat{\beta}_{3}\right) }\text{ and }%
T_{4}=\frac{\hat{\beta}_{4}}{SE\left( \hat{\beta}_{4}\right) }.
\end{equation*}%
%TCIMACRO{\TeXButton{pause}{\pause}}%
%BeginExpansion
\pause%
%EndExpansion

\item We can use $T_{3}$ and $T_{4}$ to test significance of $\beta _{3}$
and $\beta _{4}$ 
%TCIMACRO{\TeXButton{\color{red} CTRL+ cr}{\color{red}}}%
%BeginExpansion
\color{red}%
%EndExpansion
separately 
%TCIMACRO{\TeXButton{\color{black} CTRL+cb}{\color{black}}}%
%BeginExpansion
\color{black}%
%EndExpansion
by using two separate 
%TCIMACRO{\TeXButton{\color{blue} CTRL+cbl}{\color{blue}}}%
%BeginExpansion
\color{blue}%
%EndExpansion
size $\alpha $ 
%TCIMACRO{\TeXButton{\color{black} CTRL+cb}{\color{black}}}%
%BeginExpansion
\color{black}%
%EndExpansion
tests:%
%TCIMACRO{\TeXButton{pause}{\pause}}%
%BeginExpansion
\pause%
%EndExpansion

\begin{itemize}
\item Reject $H_{0,3}:\beta _{3}=0$ in favor of $H_{1,3}:\beta _{3}\neq 0$
when $\left\vert T_{3}\right\vert >t_{n-k-1,1-%
%TCIMACRO{\TeXButton{\color{red} CTRL+ cr}{\color{red}}}%
%BeginExpansion
\color{red}%
%EndExpansion
\alpha 
%TCIMACRO{\TeXButton{\color{black} CTRL+cb}{\color{black}}}%
%BeginExpansion
\color{black}%
%EndExpansion
/2}.$%
%TCIMACRO{\TeXButton{pause}{\pause}}%
%BeginExpansion
\pause%
%EndExpansion

\item Reject $H_{0,4}:\beta _{4}=0$ in favor of $H_{1,4}:\beta _{4}\neq 0$
when $\left\vert T_{4}\right\vert >t_{n-k-1,1-%
%TCIMACRO{\TeXButton{\color{red} CTRL+ cr}{\color{red}}}%
%BeginExpansion
\color{red}%
%EndExpansion
\alpha 
%TCIMACRO{\TeXButton{\color{black} CTRL+cb}{\color{black}}}%
%BeginExpansion
\color{black}%
%EndExpansion
/2}.$
\end{itemize}
\end{itemize}

%TCIMACRO{\TeXButton{EndFrame}{\end{frame}}}%
%BeginExpansion
\end{frame}%
%EndExpansion

%TCIMACRO{\TeXButton{BeginFrame}{\begin{frame}}}%
%BeginExpansion
\begin{frame}%
%EndExpansion

\QTR{frametitle}{%
%TCIMACRO{\TeXButton{$t$}{$t$}}%
%BeginExpansion
$t$%
%EndExpansion
-statistics and multiple restrictions}

\begin{itemize}
\item Rejecting $H_{0}:\beta _{3}=0,\beta _{4}=0$ in favor of $H_{1}:\beta
_{3}\neq 0$ or $\beta _{4}\neq 0$ when 
%TCIMACRO{\TeXButton{\color{blue} CTRL+cbl}{\color{blue}}}%
%BeginExpansion
\color{blue}%
%EndExpansion
at least 
%TCIMACRO{\TeXButton{\color{black} CTRL+cb}{\color{black}}}%
%BeginExpansion
\color{black}%
%EndExpansion
one of the two coefficients is significant at 
%TCIMACRO{\TeXButton{\color{blue} CTRL+cbl}{\color{blue}}}%
%BeginExpansion
\color{blue}%
%EndExpansion
level $\alpha $ 
%TCIMACRO{\TeXButton{\color{black} CTRL+cb}{\color{black}}}%
%BeginExpansion
\color{black}%
%EndExpansion
, i.e. when%
\begin{equation*}
\left\vert T_{3}\right\vert >t_{n-k-1,1-%
%TCIMACRO{\TeXButton{\color{red} CTRL+ cr}{\color{red}}}%
%BeginExpansion
\color{red}%
%EndExpansion
\alpha 
%TCIMACRO{\TeXButton{\color{black} CTRL+cb}{\color{black}}}%
%BeginExpansion
\color{black}%
%EndExpansion
/2}\text{ or }\left\vert T_{4}\right\vert >t_{n-k-1,1-%
%TCIMACRO{\TeXButton{\color{red} CTRL+ cr}{\color{red}}}%
%BeginExpansion
\color{red}%
%EndExpansion
\alpha 
%TCIMACRO{\TeXButton{\color{black} CTRL+cb}{\color{black}}}%
%BeginExpansion
\color{black}%
%EndExpansion
/2},
\end{equation*}%
%TCIMACRO{\TeXButton{\color{red} CTRL+ cr}{\color{red}}}%
%BeginExpansion
\color{red}%
%EndExpansion
is not 
%TCIMACRO{\TeXButton{\color{black} CTRL+cb}{\color{black}}}%
%BeginExpansion
\color{black}%
%EndExpansion
a 
%TCIMACRO{\TeXButton{\color{blue} CTRL+cbl}{\color{blue}}}%
%BeginExpansion
\color{blue}%
%EndExpansion
size $\alpha $ test%
%TCIMACRO{\TeXButton{\color{black} CTRL+cb}{\color{black}}}%
%BeginExpansion
\color{black}%
%EndExpansion
! 
%TCIMACRO{\TeXButton{pause}{\pause}}%
%BeginExpansion
\pause%
%EndExpansion

\begin{itemize}
\item Recall that if $A$ and $B$ are two sets 
%TCIMACRO{\TeXButton{pause}{\pause}}%
%BeginExpansion
\pause%
%EndExpansion
then $\left( A\cap B\right) \subset A$ 
%TCIMACRO{\TeXButton{pause}{\pause}}%
%BeginExpansion
\pause%
%EndExpansion
and therefore $P\left( A\cap B\right) \leq P\left( A\right) .$%
%TCIMACRO{\TeXButton{pause}{\pause}}%
%BeginExpansion
\pause%
%EndExpansion

\item 
%TCIMACRO{\TeXButton{\color{blue} CTRL+cbl}{\color{blue}}}%
%BeginExpansion
\color{blue}%
%EndExpansion
When $\beta _{3}=\beta _{4}=0$ 
%TCIMACRO{\TeXButton{\color{black} CTRL+cb}{\color{black}}}%
%BeginExpansion
\color{black}%
%EndExpansion
:%
\begin{eqnarray*}
&&P\left( \text{Reject }H_{0,3}\text{ or }H_{0,4}\right) = \\
%TCIMACRO{\TeXButton{pause}{\pause} }%
%BeginExpansion
\pause
%EndExpansion
&=&P\left( \left\vert T_{3}\right\vert >t_{n-k-1,1-%
%TCIMACRO{\TeXButton{\color{red} CTRL+ cr}{\color{red}}}%
%BeginExpansion
\color{red}%
%EndExpansion
\alpha 
%TCIMACRO{\TeXButton{\color{black} CTRL+cb}{\color{black}}}%
%BeginExpansion
\color{black}%
%EndExpansion
/2}\text{ }%
%TCIMACRO{\TeXButton{\color{blue} CTRL+cbl}{\color{blue}}}%
%BeginExpansion
\color{blue}%
%EndExpansion
\text{or }%
%TCIMACRO{\TeXButton{\color{black} CTRL+cb}{\color{black}}}%
%BeginExpansion
\color{black}%
%EndExpansion
\left\vert T_{4}\right\vert >t_{n-k-1,1-%
%TCIMACRO{\TeXButton{\color{red} CTRL+ cr}{\color{red}}}%
%BeginExpansion
\color{red}%
%EndExpansion
\alpha 
%TCIMACRO{\TeXButton{\color{black} CTRL+cb}{\color{black}}}%
%BeginExpansion
\color{black}%
%EndExpansion
/2}\right) \\
%TCIMACRO{\TeXButton{pause}{\pause} }%
%BeginExpansion
\pause
%EndExpansion
&=&P\left( \left\vert T_{3}\right\vert >t_{n-k-1,1-%
%TCIMACRO{\TeXButton{\color{red} CTRL+ cr}{\color{red}}}%
%BeginExpansion
\color{red}%
%EndExpansion
\alpha 
%TCIMACRO{\TeXButton{\color{black} CTRL+cb}{\color{black}}}%
%BeginExpansion
\color{black}%
%EndExpansion
/2}\right) +P\left( \left\vert T_{4}\right\vert >t_{n-k-1,1-%
%TCIMACRO{\TeXButton{\color{red} CTRL+ cr}{\color{red}}}%
%BeginExpansion
\color{red}%
%EndExpansion
\alpha 
%TCIMACRO{\TeXButton{\color{black} CTRL+cb}{\color{black}}}%
%BeginExpansion
\color{black}%
%EndExpansion
/2}\right) \\
&&-P\left( \left\vert T_{3}\right\vert >t_{n-k-1,1-%
%TCIMACRO{\TeXButton{\color{red} CTRL+ cr}{\color{red}}}%
%BeginExpansion
\color{red}%
%EndExpansion
\alpha 
%TCIMACRO{\TeXButton{\color{black} CTRL+cb}{\color{black}}}%
%BeginExpansion
\color{black}%
%EndExpansion
/2}\text{ }%
%TCIMACRO{\TeXButton{\color{blue} CTRL+cbl}{\color{blue}}}%
%BeginExpansion
\color{blue}%
%EndExpansion
\text{and }%
%TCIMACRO{\TeXButton{\color{black} CTRL+cb}{\color{black}}}%
%BeginExpansion
\color{black}%
%EndExpansion
\left\vert T_{4}\right\vert >t_{n-k-1,1-%
%TCIMACRO{\TeXButton{\color{red} CTRL+ cr}{\color{red}}}%
%BeginExpansion
\color{red}%
%EndExpansion
\alpha 
%TCIMACRO{\TeXButton{\color{black} CTRL+cb}{\color{black}}}%
%BeginExpansion
\color{black}%
%EndExpansion
/2}\right) \\
%TCIMACRO{\TeXButton{pause}{\pause} }%
%BeginExpansion
\pause
%EndExpansion
&=&\alpha +\alpha -P\left( \left\vert T_{3}\right\vert >t_{n-k-1,1-%
%TCIMACRO{\TeXButton{\color{red} CTRL+ cr}{\color{red}}}%
%BeginExpansion
\color{red}%
%EndExpansion
\alpha 
%TCIMACRO{\TeXButton{\color{black} CTRL+cb}{\color{black}}}%
%BeginExpansion
\color{black}%
%EndExpansion
/2}\text{ }%
%TCIMACRO{\TeXButton{\color{blue} CTRL+cbl}{\color{blue}}}%
%BeginExpansion
\color{blue}%
%EndExpansion
\text{and }%
%TCIMACRO{\TeXButton{\color{black} CTRL+cb}{\color{black}}}%
%BeginExpansion
\color{black}%
%EndExpansion
\left\vert T_{4}\right\vert >t_{n-k-1,1-%
%TCIMACRO{\TeXButton{\color{red} CTRL+ cr}{\color{red}}}%
%BeginExpansion
\color{red}%
%EndExpansion
\alpha 
%TCIMACRO{\TeXButton{\color{black} CTRL+cb}{\color{black}}}%
%BeginExpansion
\color{black}%
%EndExpansion
/2}\right) \\
%TCIMACRO{\TeXButton{pause}{\pause} }%
%BeginExpansion
\pause
%EndExpansion
&\geq &\alpha .
\end{eqnarray*}
\end{itemize}
\end{itemize}

%TCIMACRO{\TeXButton{EndFrame}{\end{frame}}}%
%BeginExpansion
\end{frame}%
%EndExpansion

%TCIMACRO{\TeXButton{BeginFrame}{\begin{frame}}}%
%BeginExpansion
\begin{frame}%
%EndExpansion

\QTR{frametitle}{Testing multiple exclusion restrictions}

\begin{itemize}
\item Consider the model%
\begin{equation*}
Y_{i}=\beta _{0}+\beta _{1}X_{1,i}+\ldots +\beta _{q}X_{q,i}+\beta
_{q+1}X_{q+1,i}+\ldots +\beta _{k}X_{k,i}+U_{i}.
\end{equation*}%
Suppose that we want to test that the first $q$ regressors have no effect on 
$Y$ (after controlling for other regressors).%
%TCIMACRO{\TeXButton{pause}{\pause}}%
%BeginExpansion
\pause%
%EndExpansion

\item The null hypothesis has $q$ 
%TCIMACRO{\TeXButton{\color{blue} CTRL+cbl}{\color{blue}}}%
%BeginExpansion
\color{blue}%
%EndExpansion
exclusion 
%TCIMACRO{\TeXButton{\color{black} CTRL+cb}{\color{black}} }%
%BeginExpansion
\color{black}
%EndExpansion
restrictions:%
\begin{equation*}
H_{0}:\beta _{1}=0,\beta _{2}=0,\ldots ,\beta _{q}=0.
\end{equation*}%
%TCIMACRO{\TeXButton{pause}{\pause}}%
%BeginExpansion
\pause%
%EndExpansion

\item The alternative hypothesis is that at least one of the restrictions in 
$H_{0}$ is false:%
\begin{equation*}
H_{1}:\beta _{1}\neq 0\text{ or }\beta _{2}\neq 0\text{ or \ldots\ or }\beta
_{q}\neq 0.
\end{equation*}
\end{itemize}

%TCIMACRO{\TeXButton{EndFrame}{\end{frame}}}%
%BeginExpansion
\end{frame}%
%EndExpansion

%TCIMACRO{\TeXButton{BeginFrame}{\begin{frame}[shrink]}}%
%BeginExpansion
\begin{frame}[shrink]%
%EndExpansion

\QTR{frametitle}{%
%TCIMACRO{\TeXButton{$F$}{$F$}}%
%BeginExpansion
$F$%
%EndExpansion
-statistic}

\begin{itemize}
\item The idea of the test is to compare the \textbf{fit} of the 
%TCIMACRO{\TeXButton{\color{blue} CTRL+cbl}{\color{blue}}}%
%BeginExpansion
\color{blue}%
%EndExpansion
unrestricted model 
%TCIMACRO{\TeXButton{\color{black} CTRL+cb}{\color{black}} }%
%BeginExpansion
\color{black}
%EndExpansion
with that of the 
%TCIMACRO{\TeXButton{\color{red} CTRL+ cr}{\color{red}}}%
%BeginExpansion
\color{red}%
%EndExpansion
null-restricted model%
%TCIMACRO{\TeXButton{\color{black} CTRL+cb}{\color{black}}}%
%BeginExpansion
\color{black}%
%EndExpansion
.%
%TCIMACRO{\TeXButton{pause}{\pause}}%
%BeginExpansion
\pause%
%EndExpansion

\item Let%
%TCIMACRO{\TeXButton{\color{blue} CTRL+cbl}{\color{blue}} }%
%BeginExpansion
\color{blue}
%EndExpansion
$SSR_{ur}$ 
%TCIMACRO{\TeXButton{\color{black} CTRL+cb}{\color{black}}}%
%BeginExpansion
\color{black}%
%EndExpansion
denote the Residual Sum-of-Squares of the 
%TCIMACRO{\TeXButton{\color{blue} CTRL+cbl}{\color{blue}}}%
%BeginExpansion
\color{blue}%
%EndExpansion
unrestricted model 
%TCIMACRO{\TeXButton{\color{black} CTRL+cb}{\color{black}}}%
%BeginExpansion
\color{black}%
%EndExpansion
\begin{equation*}
Y_{i}=\beta _{0}+\beta _{1}X_{1,i}+\ldots +\beta _{q}X_{q,i}+\beta
_{q+1}X_{q+1,i}+\ldots +\beta _{k}X_{k,i}+U_{i}.
\end{equation*}%
%TCIMACRO{\TeXButton{pause}{\pause}}%
%BeginExpansion
\pause%
%EndExpansion

\item The 
%TCIMACRO{\TeXButton{\color{red} CTRL+ cr}{\color{red}}}%
%BeginExpansion
\color{red}%
%EndExpansion
restricted model 
%TCIMACRO{\TeXButton{\color{black} CTRL+cb}{\color{black}}}%
%BeginExpansion
\color{black}%
%EndExpansion
given $H_{0}:\beta _{1}=0,\ldots ,\beta _{q}=0$ is%
\begin{equation*}
Y_{i}=\beta _{0}+\beta _{q+1}X_{q+1,i}+\ldots +\beta _{k}X_{k,i}+U_{i}.
\end{equation*}%
%TCIMACRO{\TeXButton{pause}{\pause}}%
%BeginExpansion
\pause%
%EndExpansion

\begin{itemize}
\item Let%
%TCIMACRO{\TeXButton{\color{red} CTRL+ cr}{\color{red}} }%
%BeginExpansion
\color{red}
%EndExpansion
$SSR_{r}$ 
%TCIMACRO{\TeXButton{\color{black} CTRL+cb}{\color{black}}}%
%BeginExpansion
\color{black}%
%EndExpansion
denote the Residual Sum-of-Squares of the 
%TCIMACRO{\TeXButton{\color{red} CTRL+ cr}{\color{red}}}%
%BeginExpansion
\color{red}%
%EndExpansion
restricted model 
%TCIMACRO{\TeXButton{\color{black} CTRL+cb}{\color{black}}}%
%BeginExpansion
\color{black}%
%EndExpansion
.
\end{itemize}

\item Consider the following statistic:%
\begin{equation*}
F=\frac{\left( 
%TCIMACRO{\TeXButton{\color{red} CTRL+ cr}{\color{red}}}%
%BeginExpansion
\color{red}%
%EndExpansion
SSR_{r}%
%TCIMACRO{\TeXButton{\color{black} CTRL+cb}{\color{black}}}%
%BeginExpansion
\color{black}%
%EndExpansion
-%
%TCIMACRO{\TeXButton{\color{blue} CTRL+cbl}{\color{blue}}}%
%BeginExpansion
\color{blue}%
%EndExpansion
SSR_{ur}%
%TCIMACRO{\TeXButton{\color{black} CTRL+cb}{\color{black}}}%
%BeginExpansion
\color{black}%
%EndExpansion
\right) /q}{SSR_{ur}/\left( n-k-1\right) }.
\end{equation*}%
%TCIMACRO{\TeXButton{pause}{\pause}}%
%BeginExpansion
\pause%
%EndExpansion

\begin{itemize}
\item Note that $q$\textbf{\ }=\textbf{\ number of restrictions};

\item $n-k-1$ = \textbf{unrestricted residual df, }where $k$ is the number
of regressors in the \textbf{unrestricted model.}
\end{itemize}
\end{itemize}

%TCIMACRO{\TeXButton{EndFrame}{\end{frame}}}%
%BeginExpansion
\end{frame}%
%EndExpansion

%TCIMACRO{\TeXButton{BeginFrame}{\begin{frame}}}%
%BeginExpansion
\begin{frame}%
%EndExpansion

\QTR{frametitle}{%
%TCIMACRO{\TeXButton{$F$}{$F$}}%
%BeginExpansion
$F$%
%EndExpansion
-statistic }%
\begin{equation*}
F=\frac{\left( SSR_{r}-SSR_{ur}\right) /q}{SSR_{ur}/\left( n-k-1\right) }.
\end{equation*}%
%TCIMACRO{\TeXButton{pause}{\pause}}%
%BeginExpansion
\pause%
%EndExpansion

\begin{itemize}
\item Since SSR can only increase when you drop some regressors,%
\begin{equation*}
SSR_{r}-SSR_{ur}\geq 0,
\end{equation*}%
and therefore 
%TCIMACRO{\TeXButton{\color{blue} CTRL+cbl}{\color{blue}}}%
%BeginExpansion
\color{blue}%
%EndExpansion
$F\geq 0$%
%TCIMACRO{\TeXButton{\color{black} CTRL+cb}{\color{black}}}%
%BeginExpansion
\color{black}%
%EndExpansion
.%
%TCIMACRO{\TeXButton{pause}{\pause}}%
%BeginExpansion
\pause%
%EndExpansion

\item If the null restrictions are 
%TCIMACRO{\TeXButton{\color{blue} CTRL+cbl}{\color{blue}}}%
%BeginExpansion
\color{blue}%
%EndExpansion
true%
%TCIMACRO{\TeXButton{\color{black} CTRL+cb}{\color{black}}}%
%BeginExpansion
\color{black}%
%EndExpansion
, the excluded variables do not contribute to explaining $Y$ (in
population), and therefore we should expect that $SSR_{r}-SSR_{ur}$ is small
and 
%TCIMACRO{\TeXButton{\color{blue} CTRL+cbl}{\color{blue}}}%
%BeginExpansion
\color{blue}%
%EndExpansion
$F$ is close to zero%
%TCIMACRO{\TeXButton{\color{black} CTRL+cb}{\color{black}}}%
%BeginExpansion
\color{black}%
%EndExpansion
.%
%TCIMACRO{\TeXButton{pause}{\pause}}%
%BeginExpansion
\pause%
%EndExpansion

\item If the null restriction are 
%TCIMACRO{\TeXButton{\color{red} CTRL+ cr}{\color{red}}}%
%BeginExpansion
\color{red}%
%EndExpansion
false%
%TCIMACRO{\TeXButton{\color{black} CTRL+cb}{\color{black}}}%
%BeginExpansion
\color{black}%
%EndExpansion
, the imposed restriction should substantially worsen the fit, and we should
expect that $SSR_{r}-SSR_{ur}$ is large and 
%TCIMACRO{\TeXButton{\color{red} CTRL+ cr}{\color{red}}}%
%BeginExpansion
\color{red}%
%EndExpansion
$F$ is far from zero%
%TCIMACRO{\TeXButton{\color{black} CTRL+cb}{\color{black}}}%
%BeginExpansion
\color{black}%
%EndExpansion
.%
%TCIMACRO{\TeXButton{pause}{\pause}}%
%BeginExpansion
\pause%
%EndExpansion

\item Thus, we should reject $H_{0}$ when $F>c$ where $c$ is some positive
constant.
\end{itemize}

%TCIMACRO{\TeXButton{EndFrame}{\end{frame}}}%
%BeginExpansion
\end{frame}%
%EndExpansion

%TCIMACRO{\TeXButton{BeginFrame}{\begin{frame}}}%
%BeginExpansion
\begin{frame}%
%EndExpansion

\QTR{frametitle}{%
%TCIMACRO{\TeXButton{$F$}{$F$} }%
%BeginExpansion
$F$
%EndExpansion
test}

\begin{equation*}
F=\frac{\left( SSR_{r}-SSR_{ur}\right) /q}{SSR_{ur}/\left( n-k-1\right) }.
\end{equation*}%
%TCIMACRO{\TeXButton{pause}{\pause}}%
%BeginExpansion
\pause%
%EndExpansion

\begin{itemize}
\item We should reject $H_{0}$ when $F>c.$%
%TCIMACRO{\TeXButton{pause}{\pause}}%
%BeginExpansion
\pause%
%EndExpansion

\item There is a probability that $F>c$ even when $H_{0}$ is true, thus we
need to choose $c$ so that $P\left( F>c|H_{0}\text{ is true}\right) =\alpha
. $%
%TCIMACRO{\TeXButton{pause}{\pause}}%
%BeginExpansion
\pause%
%EndExpansion

\item It turns out that when $H_{0}$ is true, the $F$-statistic has $F$
distribution with two parameters: the numerator df ($q$) and the denominator
df ($n-k-1$)%
%TCIMACRO{\TeXButton{pause}{\pause}}%
%BeginExpansion
\pause%
%EndExpansion
:%
\begin{equation*}
F\sim F_{q,n-k-1}.
\end{equation*}%
%TCIMACRO{\TeXButton{pause}{\pause}}%
%BeginExpansion
\pause%
%EndExpansion

\item Similarly to the standard normal and $t$ distributions, the $F$
distribution has been tabulated and its critical values are available in
statistical tables and statistical software such as Stata.
\end{itemize}

%TCIMACRO{\TeXButton{EndFrame}{\end{frame}}}%
%BeginExpansion
\end{frame}%
%EndExpansion

%TCIMACRO{\TeXButton{BeginFrame}{\begin{frame}}}%
%BeginExpansion
\begin{frame}%
%EndExpansion

\QTR{frametitle}{%
%TCIMACRO{\TeXButton{$F$}{$F$} }%
%BeginExpansion
$F$
%EndExpansion
test}

When $H_{0}$ is true,

\begin{equation*}
F=\frac{\left( SSR_{r}-SSR_{ur}\right) /q}{SSR_{ur}/\left( n-k-1\right) }%
\sim F_{q,n-k-1}.
\end{equation*}%
%TCIMACRO{\TeXButton{pause}{\pause}}%
%BeginExpansion
\pause%
%EndExpansion

\begin{itemize}
\item Let $F_{q,n-k-1,\tau }$ be the $\tau $-quantile of the $F_{q,n-k-1}$
distribution.%
%TCIMACRO{\TeXButton{pause}{\pause}}%
%BeginExpansion
\pause%
%EndExpansion

\item A size $\alpha $ test $H_{0}:\beta _{1}=0,\ldots ,\beta _{q}=0$
against $H_{1}:\beta _{1}\neq 0$ or \ldots\ or $\beta _{q}\neq 0$ is%
%TCIMACRO{\TeXButton{pause}{\pause}}%
%BeginExpansion
\pause%
%EndExpansion
\begin{equation*}
\text{Reject }H_{0}\text{ when }F>F_{q,n-k-1,1-\alpha }.
\end{equation*}%
%TCIMACRO{\TeXButton{pause}{\pause}}%
%BeginExpansion
\pause%
%EndExpansion

\item One can find the $p$-value by finding $\tau $ such that $%
F=F_{q,n-k-1,1-\tau }.$ The $p$-value is equal to $\tau .$
\end{itemize}

%TCIMACRO{\TeXButton{EndFrame}{\end{frame}}}%
%BeginExpansion
\end{frame}%
%EndExpansion

%TCIMACRO{\TeXButton{BeginFrame}{\begin{frame}}}%
%BeginExpansion
\begin{frame}%
%EndExpansion

\QTR{frametitle}{%
%TCIMACRO{\TeXButton{$F$}{$F$} }%
%BeginExpansion
$F$
%EndExpansion
distribution in Stata}

\begin{itemize}
\item To compute $F$ critical values use 
\begin{equation*}
\text{disp invFtail}(q,n-k-1,%
%TCIMACRO{\TeXButton{\color{red} CTRL+ cr}{\color{red}}}%
%BeginExpansion
\color{red}%
%EndExpansion
\alpha 
%TCIMACRO{\TeXButton{\color{black} CTRL+cb}{\color{black}}}%
%BeginExpansion
\color{black}%
%EndExpansion
).
\end{equation*}%
%TCIMACRO{\TeXButton{pause}{\pause}}%
%BeginExpansion
\pause%
%EndExpansion

\item To compute $p$-values from $F$ distribution use 
\begin{equation*}
\text{disp Ftail}(q,n-k-1,%
%TCIMACRO{\TeXButton{\color{red} CTRL+ cr}{\color{red}}}%
%BeginExpansion
\color{red}%
%EndExpansion
F%
%TCIMACRO{\TeXButton{\color{black} CTRL+cb}{\color{black}}}%
%BeginExpansion
\color{black}%
%EndExpansion
).
\end{equation*}
\end{itemize}

%TCIMACRO{\TeXButton{EndFrame}{\end{frame}}}%
%BeginExpansion
\end{frame}%
%EndExpansion

%TCIMACRO{\TeXButton{BeginFrame}{\begin{frame}}}%
%BeginExpansion
\begin{frame}%
%EndExpansion

\QTR{frametitle}{Example}

\begin{itemize}
\item Consider the model:%
\begin{multline*}
\ln \left( Wage_{i}\right) =\beta _{0}+\beta _{1}Experience_{i}+\beta
_{2}Experience_{i}^{2}+ \\
+\beta _{3}PrevExperience_{i}+\beta _{4}PrevExperience_{i}^{2}+\beta
_{5}Education_{i}+U_{i}.
\end{multline*}

\item We test%
\begin{equation*}
H_{0}:\beta _{3}=0,\beta _{4}=0\text{ against }H_{1}:\beta _{3}\neq 0\text{
or }\beta _{4}\neq 0.
\end{equation*}%
%TCIMACRO{\TeXButton{pause}{\pause}}%
%BeginExpansion
\pause%
%EndExpansion

\item $q=2.$%
%TCIMACRO{\TeXButton{pause}{\pause}}%
%BeginExpansion
\pause%
%EndExpansion

\item $\alpha =0.05.$
\end{itemize}

%TCIMACRO{\TeXButton{EndFrame}{\end{frame}}}%
%BeginExpansion
\end{frame}%
%EndExpansion

%TCIMACRO{%
%\TeXButton{BeginFrame}{\begin{frame}[fragile=singleslide,shrink=30]}}%
%BeginExpansion
\begin{frame}[fragile=singleslide,shrink=30]%
%EndExpansion

\QTR{frametitle}{Example: the unrestricted model}
\begin{verbatim}
. regress lnWage Experience Experience2 PrevExperience PrevExperience2 Education
 
      Source |       SS       df       MS              Number of obs =     526
-------------+------------------------------           F(  5,   520) =   55.04
       Model |  51.3318741     5  10.2663748           Prob > F      =  0.0000
    Residual |  96.9978773   520  .186534379           R-squared     =  0.3461
-------------+------------------------------           Adj R-squared =  0.3398
       Total |  148.329751   525   .28253286           Root MSE      =   .4319
 
------------------------------------------------------------------------------
      lnWage |      Coef.   Std. Err.      t    P>|t|     [95% Conf. Interval]
-------------+----------------------------------------------------------------
  Experience |   .0471914   .0068074     6.93   0.000     .0338179    .0605649
 Experience2 |  -.0008518   .0002472    -3.45   0.001    -.0013374   -.0003662
PrevExperi~e |   .0168997   .0047331     3.57   0.000     .0076013    .0261981
PrevExperi~2 |  -.0003727   .0001208    -3.09   0.002      -.00061   -.0001354
   Education |   .0887704   .0072131    12.31   0.000     .0745999    .1029408
       _cons |   .2368427     .10287     2.30   0.022     .0347509    .4389346
------------------------------------------------------------------------------
\end{verbatim}

\begin{itemize}
\item $SSR_{ur}=$96.9978773.

\item $n-k-1=$526-5-1=520.
\end{itemize}

%TCIMACRO{\TeXButton{EndFrame}{\end{frame}}}%
%BeginExpansion
\end{frame}%
%EndExpansion

%TCIMACRO{%
%\TeXButton{BeginFrame}{\begin{frame}[fragile=singleslide,shrink=30]}}%
%BeginExpansion
\begin{frame}[fragile=singleslide,shrink=30]%
%EndExpansion

\QTR{frametitle}{Example: the restricted model}
\begin{verbatim}
. regress lnWage Experience Experience2  Education
 
      Source |       SS       df       MS              Number of obs =     526
-------------+------------------------------           F(  3,   522) =   85.49
       Model |  48.8668114     3  16.2889371           Prob > F      =  0.0000
    Residual |    99.46294   522  .190542031           R-squared     =  0.3294
-------------+------------------------------           Adj R-squared =  0.3256
       Total |  148.329751   525   .28253286           Root MSE      =  .43651
 
------------------------------------------------------------------------------
      lnWage |      Coef.   Std. Err.      t    P>|t|     [95% Conf. Interval]
-------------+----------------------------------------------------------------
  Experience |   .0510784   .0067937     7.52   0.000      .037732    .0644248
 Experience2 |  -.0009941   .0002463    -4.04   0.000     -.001478   -.0005103
   Education |   .0852822   .0068978    12.36   0.000     .0717313    .0988331
       _cons |   .3688491   .0908138     4.06   0.000     .1904437    .5472544
------------------------------------------------------------------------------
\end{verbatim}

\begin{itemize}
\item $SSR_{r}=$99.46294.
\end{itemize}

%TCIMACRO{\TeXButton{EndFrame}{\end{frame}}}%
%BeginExpansion
\end{frame}%
%EndExpansion

%TCIMACRO{\TeXButton{BeginFrame}{\begin{frame}[squeeze]}}%
%BeginExpansion
\begin{frame}[squeeze]%
%EndExpansion

\QTR{frametitle}{Example: 
%TCIMACRO{\TeXButton{$F$}{$F$} }%
%BeginExpansion
$F$
%EndExpansion
statistic and test}

\begin{itemize}
\item To compute the statistic:%
\begin{equation*}
F=\frac{\left( SSR_{r}-SSR_{ur}\right) /q}{SSR_{ur}/\left( n-k-1\right) }=%
\frac{\left( 99.46294-96.9978773\right) /2}{96.9978773/\left( 526-5-1\right) 
}\approx 6.61.
\end{equation*}%
%TCIMACRO{\TeXButton{pause}{\pause}}%
%BeginExpansion
\pause%
%EndExpansion

\item The critical value:

%TCIMACRO{\TeXButton{\begin{semiverbatim}}{\begin{semiverbatim}}}%
%BeginExpansion
\begin{semiverbatim}%
%EndExpansion
.\ disp\ invFtail(2,520,0.05)

3.0130572%
%TCIMACRO{\TeXButton{\end{semiverbatim}}{\end{semiverbatim}}}%
%BeginExpansion
\end{semiverbatim}%
%EndExpansion
%TCIMACRO{\TeXButton{pause}{\pause}}%
%BeginExpansion
\pause%
%EndExpansion

\item The test: $6.61>$3.0130572 and at 5\% significance level we reject $%
H_{0}$ that previous experience has no effect on wage.%
%TCIMACRO{\TeXButton{pause}{\pause}}%
%BeginExpansion
\pause%
%EndExpansion

\item The $p$-value:

%TCIMACRO{\TeXButton{\begin{semiverbatim}}{\begin{semiverbatim}}}%
%BeginExpansion
\begin{semiverbatim}%
%EndExpansion
.\ disp\ Ftail(2,520,6.61)

.00146284%
%TCIMACRO{\TeXButton{\end{semiverbatim}}{\end{semiverbatim}}}%
%BeginExpansion
\end{semiverbatim}%
%EndExpansion
%TCIMACRO{\TeXButton{pause}{\pause}}%
%BeginExpansion
\pause%
%EndExpansion

$\Longrightarrow $ We reject $H_{0}$ for any $\alpha >0.00146284.$
\end{itemize}

%TCIMACRO{\TeXButton{EndFrame}{\end{frame}}}%
%BeginExpansion
\end{frame}%
%EndExpansion

%TCIMACRO{%
%\TeXButton{BeginFrame}{\begin{frame}
%}}%
%BeginExpansion
\begin{frame}
%
%EndExpansion

\QTR{frametitle}{Example: Stata test command}

\begin{itemize}
\item Instead of running two models, restricted and unrestricted, one can
use the Stata 
%TCIMACRO{\TeXButton{\color{blue} CTRL+cbl}{\color{blue}}}%
%BeginExpansion
\color{blue}%
%EndExpansion
test 
%TCIMACRO{\TeXButton{\color{black} CTRL+cb}{\color{black}}}%
%BeginExpansion
\color{black}%
%EndExpansion
command after estimation of the 
%TCIMACRO{\TeXButton{\color{red} CTRL+ cr}{\color{red}}}%
%BeginExpansion
\color{red}%
%EndExpansion
unrestricted model%
%TCIMACRO{\TeXButton{\color{black} CTRL+cb}{\color{black}}}%
%BeginExpansion
\color{black}%
%EndExpansion
.%
%TCIMACRO{\TeXButton{pause}{\pause}}%
%BeginExpansion
\pause%
%EndExpansion

\item To test that previous experience has no effect:\medskip 
%TCIMACRO{\TeXButton{\begin{semiverbatim}}{\begin{semiverbatim}}}%
%BeginExpansion
\begin{semiverbatim}%
%EndExpansion

. test (PrevExperience=0) (PrevExperience2=0)%
%TCIMACRO{\TeXButton{\end{semiverbatim}}{\end{semiverbatim}}}%
%BeginExpansion
\end{semiverbatim}%
%EndExpansion
%TCIMACRO{\TeXButton{pause}{\pause}}%
%BeginExpansion
\pause%
%EndExpansion

\item The output of this command is:\medskip 
%TCIMACRO{\TeXButton{\begin{semiverbatim}}{\begin{semiverbatim}}}%
%BeginExpansion
\begin{semiverbatim}%
%EndExpansion

( 1) PrevExperience = 0

( 2) PrevExperience2 = 0

F( 2, 520) = 6.61

Prob \TEXTsymbol{>} F = 0.0015

%TCIMACRO{\TeXButton{\end{semiverbatim}}{\end{semiverbatim}}}%
%BeginExpansion
\end{semiverbatim}%
%EndExpansion
\end{itemize}

%TCIMACRO{\TeXButton{EndFrame}{\end{frame}}}%
%BeginExpansion
\end{frame}%
%EndExpansion

%TCIMACRO{\TeXButton{BeginFrame}{\begin{frame}[fragile=singleslide]}}%
%BeginExpansion
\begin{frame}[fragile=singleslide]%
%EndExpansion

\QTR{frametitle}{Example: Stata test command}

\begin{itemize}
\item To test that the coefficient on previous experience equal to the
coefficient on experience and the coefficient on previous experience squared
is zero:\medskip

.\ test\ (Experience==PrevExperience2)\ (PrevExperience2=0)

\medskip

\item The output is:\medskip

\ (\ 1)\ \ Experience\ -\ PrevExperience2\ =\ 0

\ (\ 2)\ \ PrevExperience2\ =\ 0

\ \ \ \ \ \ \ F(\ \ 2,\ \ \ 520)\ =\ \ \ 31.94

\ \ \ \ \ \ \ \ \ \ \ \ Prob\ \TEXTsymbol{>}\ F\ =\ \ \ \ 0.0000
\end{itemize}

%TCIMACRO{\TeXButton{EndFrame}{\end{frame}}}%
%BeginExpansion
\end{frame}%
%EndExpansion

%TCIMACRO{\TeXButton{BeginFrame}{\begin{frame}}}%
%BeginExpansion
\begin{frame}%
%EndExpansion

\QTR{frametitle}{%
%TCIMACRO{\TeXButton{$F$}{$F$} }%
%BeginExpansion
$F$
%EndExpansion
and 
%TCIMACRO{\TeXButton{$R^2$}{$R^2$}}%
%BeginExpansion
$R^2$%
%EndExpansion
}

\begin{itemize}
\item Let $R_{ur}^{2}$ denote the $R^{2}$ corresponding to the unrestricted
model:%
\begin{equation*}
Y_{i}=\beta _{0}+\beta _{1}X_{1,i}+\ldots +\beta _{q}X_{q,i}+\beta
_{q+1}X_{q+1,i}+\ldots +\beta _{k}X_{k,i}+U_{i}.
\end{equation*}%
%TCIMACRO{\TeXButton{pause}{\pause}}%
%BeginExpansion
\pause%
%EndExpansion

\item Let $R_{r}^{2}$ denote the $R^{2}$ corresponding to the restricted
model:%
\begin{equation*}
Y_{i}=\beta _{0}+\beta _{q+1}X_{q+1,i}+\ldots +\beta _{k}X_{k,i}+U_{i}.
\end{equation*}%
%TCIMACRO{\TeXButton{pause}{\pause}}%
%BeginExpansion
\pause%
%EndExpansion

\item The two models have the same dependent variable and therefore the same
Total Sum-of-Squares:%
\begin{equation*}
SST=\sum_{i=1}^{n}\left( Y_{i}-\bar{Y}\right) ^{2}=SST_{ur}=SST_{r}.
\end{equation*}
\end{itemize}

%TCIMACRO{\TeXButton{EndFrame}{\end{frame}}}%
%BeginExpansion
\end{frame}%
%EndExpansion

%TCIMACRO{\TeXButton{BeginFrame}{\begin{frame}}}%
%BeginExpansion
\begin{frame}%
%EndExpansion

\QTR{frametitle}{%
%TCIMACRO{\TeXButton{$F$}{$F$} }%
%BeginExpansion
$F$
%EndExpansion
and 
%TCIMACRO{\TeXButton{$R^2$}{$R^2$}}%
%BeginExpansion
$R^2$%
%EndExpansion
}

\begin{itemize}
\item In this case, we can write then%
\begin{eqnarray*}
F &=&\frac{\left( SSR_{r}-SSR_{ur}\right) /q}{SSR_{ur}/\left( n-k-1\right) }
\\
%TCIMACRO{\TeXButton{pause}{\pause} }%
%BeginExpansion
\pause
%EndExpansion
&=&\frac{\left( \frac{SSR_{r}}{SST}-\frac{SSR_{ur}}{SST}\right) /q}{\frac{%
SSR_{ur}}{SST}/\left( n-k-1\right) } \\
%TCIMACRO{\TeXButton{pause}{\pause} }%
%BeginExpansion
\pause
%EndExpansion
&=&\frac{\left( 1-R_{r}^{2}-\left( 1-R_{ur}^{2}\right) \right) /q}{\left(
1-R_{ur}^{2}\right) /\left( n-k-1\right) } \\
%TCIMACRO{\TeXButton{pause}{\pause} }%
%BeginExpansion
\pause
%EndExpansion
&=&%
%TCIMACRO{\TeXButton{\color{blue} CTRL+cbl}{\color{blue}}}%
%BeginExpansion
\color{blue}%
%EndExpansion
\frac{\left( R_{ur}^{2}-R_{r}^{2}\right) /q}{\left( 1-R_{ur}^{2}\right)
/\left( n-k-1\right) }%
%TCIMACRO{\TeXButton{\color{black} CTRL+cb}{\color{black}}}%
%BeginExpansion
\color{black}%
%EndExpansion
.
\end{eqnarray*}
\end{itemize}

%TCIMACRO{\TeXButton{EndFrame}{\end{frame}}}%
%BeginExpansion
\end{frame}%
%EndExpansion

%TCIMACRO{\TeXButton{BeginFrame}{\begin{frame}}}%
%BeginExpansion
\begin{frame}%
%EndExpansion

\QTR{frametitle}{%
%TCIMACRO{\TeXButton{$F$}{$F$} }%
%BeginExpansion
$F$
%EndExpansion
test: more examples}

\begin{itemize}
\item Suppose that you want to test $H_{0}:\beta _{1}=1$ against $%
H_{1}:\beta _{1}\neq 1$ in%
\begin{equation*}
Y_{i}=\beta _{0}+\beta _{1}X_{1,i}+\beta _{2}X_{2,i}+\ldots +\beta
_{k}X_{k,i}+U_{i}.
\end{equation*}%
%TCIMACRO{\TeXButton{pause}{\pause}}%
%BeginExpansion
\pause%
%EndExpansion

\item The 
%TCIMACRO{\TeXButton{\color{red} CTRL+ cr}{\color{red}}}%
%BeginExpansion
\color{red}%
%EndExpansion
restricted 
%TCIMACRO{\TeXButton{\color{black} CTRL+cb}{\color{black}}}%
%BeginExpansion
\color{black}%
%EndExpansion
model is%
\begin{equation*}
Y_{i}=\beta _{0}+X_{1,i}+\beta _{2}X_{2,i}+\ldots +\beta _{k}X_{k,i}+U_{i},
\end{equation*}%
%TCIMACRO{\TeXButton{pause}{\pause}}%
%BeginExpansion
\pause%
%EndExpansion
or%
\begin{equation*}
Y_{i}-X_{1,i}=\beta _{0}+\beta _{2}X_{2,i}+\ldots +\beta _{k}X_{k,i}+U_{i}.
\end{equation*}%
%TCIMACRO{\TeXButton{pause}{\pause}}%
%BeginExpansion
\pause%
%EndExpansion

\begin{enumerate}
\item Generate a new dependent variable $Y_{i}^{\ast }=Y_{i}-X_{1,i}.$%
%TCIMACRO{\TeXButton{pause}{\pause}}%
%BeginExpansion
\pause%
%EndExpansion

\item Regress $Y^{\ast }$ against a constant, $X_{2},\ldots ,X_{k}$ to
obtain $SSR_{r}.$%
%TCIMACRO{\TeXButton{pause}{\pause}}%
%BeginExpansion
\pause%
%EndExpansion

\item Estimate the unrestricted model to obtain $SSR_{ur}.$%
%TCIMACRO{\TeXButton{pause}{\pause}}%
%BeginExpansion
\pause%
%EndExpansion

\item Compute $F=\frac{\left( SSR_{r}-SSR_{ur}\right) /1}{SSR_{ur}/\left(
n-k-1\right) }.$
\end{enumerate}
\end{itemize}

%TCIMACRO{\TeXButton{EndFrame}{\end{frame}}}%
%BeginExpansion
\end{frame}%
%EndExpansion

%TCIMACRO{\TeXButton{BeginFrame}{\begin{frame}}}%
%BeginExpansion
\begin{frame}%
%EndExpansion

\QTR{frametitle}{%
%TCIMACRO{\TeXButton{$F$}{$F$} }%
%BeginExpansion
$F$
%EndExpansion
test: more examples}

\begin{itemize}
\item Suppose that you want to test $H_{0}:\beta _{1}+\beta _{2}=1$ against $%
H_{1}:\beta _{1}+\beta _{2}\neq 1$ in%
\begin{equation*}
Y_{i}=\beta _{0}+\beta _{1}X_{1,i}+\beta _{2}X_{2,i}+\ldots +\beta
_{k}X_{k,i}+U_{i}.
\end{equation*}%
%TCIMACRO{\TeXButton{pause}{\pause}}%
%BeginExpansion
\pause%
%EndExpansion

\item The 
%TCIMACRO{\TeXButton{\color{red} CTRL+ cr}{\color{red}}}%
%BeginExpansion
\color{red}%
%EndExpansion
restricted 
%TCIMACRO{\TeXButton{\color{black} CTRL+cb}{\color{black}}}%
%BeginExpansion
\color{black}%
%EndExpansion
model is%
\begin{equation*}
Y_{i}=\beta _{0}+\left( 1-\beta _{2}\right) X_{1,i}+\beta _{2}X_{2,i}+\ldots
+\beta _{k}X_{k,i}+U_{i},
\end{equation*}%
%TCIMACRO{\TeXButton{pause}{\pause}}%
%BeginExpansion
\pause%
%EndExpansion
or%
\begin{equation*}
Y_{i}-X_{1,i}=\beta _{0}+\beta _{2}\left( X_{2,i}-X_{1,i}\right) +\ldots
+\beta _{k}X_{k,i}+U_{i}.
\end{equation*}%
%TCIMACRO{\TeXButton{pause}{\pause}}%
%BeginExpansion
\pause%
%EndExpansion

\begin{enumerate}
\item Generate a new dependent variable $Y_{i}^{\ast }=Y_{i}-X_{1,i}.$%
%TCIMACRO{\TeXButton{pause}{\pause}}%
%BeginExpansion
\pause%
%EndExpansion

\item Generate a new regressor $X_{2}^{\ast }=X_{2,i}-X_{1,i}.$%
%TCIMACRO{\TeXButton{pause}{\pause}}%
%BeginExpansion
\pause%
%EndExpansion

\item Regress $Y^{\ast }$ against a constant, $X_{2}^{\ast },X_{3},\ldots
,X_{k}$ to obtain $SSR_{r}.$%
%TCIMACRO{\TeXButton{pause}{\pause}}%
%BeginExpansion
\pause%
%EndExpansion

\item Estimate the unrestricted model to obtain $SSR_{ur}.$%
%TCIMACRO{\TeXButton{pause}{\pause}}%
%BeginExpansion
\pause%
%EndExpansion

\item Compute $F=\frac{\left( SSR_{r}-SSR_{ur}\right) /1}{SSR_{ur}/\left(
n-k-1\right) }.$
\end{enumerate}
\end{itemize}

%TCIMACRO{\TeXButton{EndFrame}{\end{frame}}}%
%BeginExpansion
\end{frame}%
%EndExpansion

%TCIMACRO{\TeXButton{BeginFrame}{\begin{frame}}}%
%BeginExpansion
\begin{frame}%
%EndExpansion

\QTR{frametitle}{Relationship between 
%TCIMACRO{\TeXButton{$F$}{$F$} }%
%BeginExpansion
$F$
%EndExpansion
and 
%TCIMACRO{\TeXButton{$t$}{$t$} }%
%BeginExpansion
$t$
%EndExpansion
statistics}

\begin{itemize}
\item The $F$ statistic can also be used for testing a single restriction.%
%TCIMACRO{\TeXButton{pause}{\pause}}%
%BeginExpansion
\pause%
%EndExpansion

\item In the case of a single restriction, the $F$ test and $t$ test lead to
the 
%TCIMACRO{\TeXButton{\color{blue} CTRL+cbl}{\color{blue}}}%
%BeginExpansion
\color{blue}%
%EndExpansion
same outcome 
%TCIMACRO{\TeXButton{\color{black} CTRL+cb}{\color{black}}}%
%BeginExpansion
\color{black}%
%EndExpansion
because 
\begin{equation*}
t_{n-k-1}^{2}=F_{1,n-k-1}.
\end{equation*}
\end{itemize}

%TCIMACRO{\TeXButton{EndFrame}{\end{frame}}}%
%BeginExpansion
\end{frame}%
%EndExpansion

%TCIMACRO{\TeXButton{BeginFrame}{\begin{frame}}}%
%BeginExpansion
\begin{frame}%
%EndExpansion

\QTR{frametitle}{Test of model significance}

\begin{itemize}
\item Consider the model%
\begin{equation*}
Y_{i}=\beta _{0}+\beta _{1}X_{1,i}+\ldots +\beta _{k}X_{k,i}+U_{i}.
\end{equation*}%
%TCIMACRO{\TeXButton{pause}{\pause}}%
%BeginExpansion
\pause%
%EndExpansion

\item Suppose that you want to test that none of the regressors explain $Y:$%
\begin{eqnarray*}
H_{0} &:&\beta _{1}=\beta _{2}=\ldots =\beta _{k}=0\text{ (}k\text{
restrictions) against } \\
H_{1} &:&\beta _{j}\neq 0\text{ for some }j=1,\ldots ,k.
\end{eqnarray*}%
%TCIMACRO{\TeXButton{pause}{\pause}}%
%BeginExpansion
\pause%
%EndExpansion

\item The restricted model is given by%
\begin{equation*}
Y_{i}=\beta _{0}+U_{i},
\end{equation*}%
and since $\hat{\beta}_{0}=\bar{Y}$ in this model,%
\begin{equation*}
SSR_{r}=\sum_{i=1}^{n}\left( Y_{i}-\bar{Y}\right) ^{2}=SST\text{ and }%
SSR_{ur}=SSR.
\end{equation*}
\end{itemize}

%TCIMACRO{\TeXButton{EndFrame}{\end{frame}}}%
%BeginExpansion
\end{frame}%
%EndExpansion

%TCIMACRO{\TeXButton{BeginFrame}{\begin{frame}}}%
%BeginExpansion
\begin{frame}%
%EndExpansion

\QTR{frametitle}{Test of model significance}

\begin{itemize}
\item The $F$ statistic for model significance test is%
\begin{eqnarray*}
F &=&\frac{\left( SSR_{r}-SSR_{ur}\right) /k}{SSR_{ur}/\left( n-k-1\right) }
\\
%TCIMACRO{\TeXButton{pause}{\pause} }%
%BeginExpansion
\pause
%EndExpansion
&=&\frac{\left( SST-SSR\right) /k}{SSR/\left( n-k-1\right) } \\
%TCIMACRO{\TeXButton{pause}{\pause} }%
%BeginExpansion
\pause
%EndExpansion
&=&\frac{SSE/k}{SSR/\left( n-k-1\right) } \\
%TCIMACRO{\TeXButton{pause}{\pause} }%
%BeginExpansion
\pause
%EndExpansion
&=&\frac{R^{2}/k}{\left( 1-R^{2}\right) /\left( n-k-1\right) }.
\end{eqnarray*}
\end{itemize}

%TCIMACRO{\TeXButton{EndFrame}{\end{frame}}}%
%BeginExpansion
\end{frame}%
%EndExpansion

%TCIMACRO{%
%\TeXButton{BeginFrame}{\begin{frame}[fragile=singleslide,shrink=30]}}%
%BeginExpansion
\begin{frame}[fragile=singleslide,shrink=30]%
%EndExpansion

\QTR{frametitle}{Test of model significance}

\begin{itemize}
\item The $F$ statistic for the model significance test and its $p$-value is
reported by Stata as in the top part of the regression output.\medskip
\end{itemize}

\begin{verbatim}
 
      Source |       SS       df       MS              Number of obs =     526
-------------+------------------------------           F(  5,   520) =   55.04
       Model |  51.3318741     5  10.2663748           Prob > F      =  0.0000
    Residual |  96.9978773   520  .186534379           R-squared     =  0.3461
-------------+------------------------------           Adj R-squared =  0.3398
       Total |  148.329751   525   .28253286           Root MSE      =   .4319
 
\end{verbatim}

%TCIMACRO{\TeXButton{EndFrame}{\end{frame}}}%
%BeginExpansion
\end{frame}%
%EndExpansion

\end{document}
